\section{How the Final Controller Was Developed}\label{sec:development}

The controller was developed in four steps: a quality-oriented profile, an
identity-aware profile, transfer and cross-pipeline studies, and a modern
stage that produced the frozen controller. Each step is compared with the
baseline it was built on. The historical steps used other detectors and
evaluation protocols than the modern stage, so their numbers are reported in
their own tables and are not pooled with the modern results.

\subsection{Quality-oriented profile}\label{sec:quality}

\emph{Baseline.} The reference pipeline is the pretrained YOLOv8n detector
\cite{jocher2023yolov8} followed by ByteTrack \cite{zhang2022bytetrack}, both at
their default settings: confidence threshold 0.25, NMS IoU threshold 0.45,
input size 640 pixels, and the default ByteTrack parameters (high threshold
0.25, low threshold 0.10, new-track threshold 0.25, buffer 30 frames, match
threshold 0.80). It runs with one fixed operating point for the whole video.

\emph{Method.} The first optimization-derived profile kept the same detector,
used a tuned ByteTrack profile, and searched only on the VisDrone2019-MOT
validation split. Five cues were used: crowd count, the fraction of earlier
boxes below $32\times32$ pixels, edge density, low brightness, and low
Laplacian variance. Each cue was normalized with validation statistics, and
the index was their weighted sum
\begin{equation}
S_j=\sum_{k\in\{c,t,e,n,b\}}w_k\,x_{k,j},
\label{eq:hist-sci}
\end{equation}
with $w_k\geq0$ and $\sum_k w_k=1$, averaged over seven analyses taken every
ten frames. Confidence and NMS were interpolated linearly between the values
at $S=0$ and $S=1$, and two sorted thresholds selected one of three
resolutions (512, 928, 960 pixels).

An Optuna TPE study \cite{akiba2019optuna} jointly searched the five weights,
the confidence and NMS end points, and the two thresholds. The resolution
levels, the supported confidence values $\{0.25,\dots,0.45\}$, the supported
NMS values from 0.30 to 0.70, and the temporal settings were fixed beforehand
by validation sweeps (Sect.~\ref{sec:constants}). The selection rule kept
trials with at least 25 FPS and at most 271 IDS (a reference value from
earlier validation runs) and then took the highest MOTA, with lower IDS, HOTA,
IDF1, and FPS as tie-breaks. The frozen profile (archived as Trial~24) has
$w_c=0.1294928$, $w_t=0.2217477$, $w_e=0.4337134$, $w_n=0.0535577$,
$w_b=0.1614885$, confidence end points $(0.30,0.40)$, NMS end points
$(0.35,0.35)$, and thresholds $(0.1353494, 0.2872868)$. On validation it
reached MOTA 23.038, HOTA 36.110, IDF1 40.758, and IDS 270 at 37.169 FPS.

\emph{Comparison.} Table~\ref{tab:quality} compares the profile with the
baseline on the historical VisDrone test-dev split, with paired bootstrap
intervals (5000 resamples, seed 42). HOTA rises by 5.4051 points, MOTA by
7.2195 points, and IDF1 by 8.8220 points, and throughput rises from 36.528 to
38.985 FPS.

\begin{table}[t]
\caption{Quality-oriented profile versus its baseline (YOLOv8n with ByteTrack
at default settings) on VisDrone test-dev, historical protocol. Gain: paired
bootstrap difference with 95\% interval; for IDS it is the reduction.}
\label{tab:quality}
\centering\footnotesize
\setlength{\tabcolsep}{4pt}
\begin{tabular}{@{}lrrl@{}}
\toprule
Metric & Baseline & Quality & Gain [95\% CI]\\
\midrule
HOTA & 28.430 & 33.835 & $+5.4051$ $[4.1273, 6.9022]$\\
MOTA & 19.729 & 26.948 & $+7.2195$ $[5.4362, 9.2934]$\\
IDF1 & 32.724 & 41.546 & $+8.8220$ $[6.9005, 11.0537]$\\
IDS & 1235 & 1184 & $51$ $[-114, 219]$\\
FPS & 36.528 & 38.985 & --\\
\bottomrule
\end{tabular}
\end{table}

\subsection{Identity-aware Pareto profile}\label{sec:identity}

\emph{Baseline.} The baseline is the same YOLOv8n and ByteTrack pipeline at
default settings (Sect.~\ref{sec:quality}).

\emph{Method.} Quality is not the only deployment goal; identity stability and
compute can move in a different direction from MOTA. The second study kept
the same cues, operating levels, temporal settings, and validation data, but
changed the question. It maximized MOTA and minimized IDS subject only to
$\mathrm{FPS}\geq25$, without the IDS gate. A seeded TPE sampler ran 49
completed trials. A trial dominated another when it had no lower MOTA and no
higher IDS, with at least one strict improvement, and the feasible Pareto
front was built from these relations \cite{deb2002nsga2}. The balanced
candidate maximized an equal-weight score of normalized MOTA and normalized IDS
reduction, with HOTA, IDF1, FPS, and trial index as tie-breaks.

The selected profile (archived as Trial~22) has weights $w_c=0.1646453$,
$w_t=0.1746265$, $w_e=0.5076113$, $w_n=0.1206956$, $w_b=0.0324213$,
confidence end points $(0.40,0.40)$, NMS end points $(0.60,0.35)$, and
thresholds $(0.2927136, 0.6661672)$. On validation it reached MOTA 19.330,
HOTA 31.651, IDF1 34.226, and IDS 168 at 51.406 FPS.

\emph{Comparison.} Table~\ref{tab:identity} compares the profile with the
baseline on test-dev. It removes 316 identity switches and raises throughput
from 36.528 to 46.024 FPS, while HOTA, MOTA, and IDF1 also rise. Against the
quality-oriented profile it has 265 fewer IDS $[82, 502]$ and $-2.617$ HOTA
$[-4.004, -1.694]$. The two profiles therefore answer different questions: one
favors overall quality and the other favors identity stability and speed.

\begin{table}[t]
\caption{Identity-aware profile versus its baseline (YOLOv8n with ByteTrack at
default settings) on VisDrone test-dev, historical protocol. Gain: paired
bootstrap difference with 95\% interval; for IDS it is the reduction.}
\label{tab:identity}
\centering\footnotesize
\setlength{\tabcolsep}{4pt}
\begin{tabular}{@{}lrrl@{}}
\toprule
Metric & Baseline & Identity & Gain [95\% CI]\\
\midrule
HOTA & 28.430 & 31.218 & $+2.788$ $[0.747, 4.660]$\\
MOTA & 19.729 & 23.792 & $+4.063$ $[0.952, 6.886]$\\
IDF1 & 32.724 & 37.870 & $+5.146$ $[2.028, 8.080]$\\
IDS & 1235 & 919 & $316$ $[175, 470]$\\
FPS & 36.528 & 46.024 & --\\
\bottomrule
\end{tabular}
\end{table}

\subsection{Transfer and cross-pipeline studies}\label{sec:transfer}

\emph{UAVDT.} The baseline is again YOLOv8n with ByteTrack at default settings
(Sect.~\ref{sec:quality}). Both frozen profiles were applied to UAVDT
\cite{du2018uavdt} without retuning (Table~\ref{tab:hist-uavdt}). The
quality-oriented profile raised HOTA by 4.305 points $[2.740, 5.709]$ and the
identity-aware profile by 2.845 points $[1.238, 4.348]$; their IDS reductions
are 237 $[104, 377]$ and 250 $[110, 404]$.

\begin{table}[t]
\caption{Historical UAVDT transfer without retuning. Baseline: YOLOv8n with
ByteTrack at default settings.}
\label{tab:hist-uavdt}
\centering\footnotesize
\setlength{\tabcolsep}{4pt}
\begin{tabular}{@{}lrrr@{}}
\toprule
Metric & Baseline & Quality & Identity\\
\midrule
HOTA & 24.085 & 28.390 & 26.930\\
MOTA & 13.841 & 17.399 & 16.118\\
IDF1 & 27.887 & 34.453 & 32.014\\
IDS & 558 & 321 & 308\\
FPS & 65.015 & 58.353 & 61.462\\
\bottomrule
\end{tabular}
\end{table}

\emph{U2MOT on VisDrone.} The baseline is the published U2MOT tracker
\cite{liu2023u2mot} with its released YOLOX-X detector \cite{ge2021yolox}, run
at its own static setting. Seven VisDrone validation sequences served for
development and 17 test-dev sequences for evaluation. The cue weights and
transforms of the quality-oriented profile were kept, the tracker was not
modified, and the action boundaries were set to the 33rd and 67th percentiles
of the index on the U2MOT validation run, without using ground truth.
Table~\ref{tab:u2mot} gives the comparison: the controlled pipeline has 87
fewer IDS and 1230 fewer FP.

\begin{table}[t]
\caption{Cross-pipeline study on VisDrone test-dev. Baseline: U2MOT with its
YOLOX-X detector at its static setting.}
\label{tab:u2mot}
\centering\footnotesize
\setlength{\tabcolsep}{4pt}
\begin{tabular}{@{}lrr@{}}
\toprule
Metric & U2MOT baseline & With controller\\
\midrule
MOTA & 53.873 & 53.767\\
IDF1 & 69.775 & 69.849\\
IDS & 1239 & 1152\\
FP & 41385 & 40155\\
FN & 63241 & 64801\\
\bottomrule
\end{tabular}
\end{table}

\emph{SparseTrack on MOT17.} The baseline is the published SparseTrack tracker
\cite{liu2025sparsetrack} with a fixed NMS IoU threshold of 0.75, on the seven
MOT17 \texttt{val\_half} sequences (2652 frames) \cite{milan2020motchallenge}.
A minimal edge-driven rule switched only the NMS threshold, between 0.70 and
0.80; its threshold came from a leave-one-sequence-out study on the same
sequences. Table~\ref{tab:sparsetrack} gives the comparison: the adaptive rule
reduces IDS from 136 to 123, with MOTA $+0.0371$ points $[-0.041, +0.182]$.

\begin{table}[t]
\caption{Cross-pipeline study on MOT17 \texttt{val\_half}. Baseline:
SparseTrack with a fixed NMS IoU threshold of 0.75.}
\label{tab:sparsetrack}
\centering\footnotesize
\setlength{\tabcolsep}{4pt}
\begin{tabular}{@{}lrr@{}}
\toprule
Metric & SparseTrack baseline & With adaptive NMS\\
\midrule
MOTA & 76.8881 & 76.9252\\
IDF1 & 81.4943 & 81.5299\\
IDS & 136 & 123\\
\bottomrule
\end{tabular}
\end{table}

Together these studies show that the preferred profile depends on the
objective (quality or identity) and on the tracker host. This motivated a
modern stage with a stronger detector, explicit calibration data, and two
tracker hosts.

\subsection{Modern development of the frozen controller}\label{sec:modern}

\emph{Baseline.} The modern baseline is the YOLO11m detector
\cite{jocher2024yolo11}, in the public VisDrone checkpoint described in
Sect.~\ref{sec:setup}, run at a fixed operating point of 1536 pixels and NMS
IoU 0.70 with no extra score threshold. It feeds ByteTrack
\cite{zhang2022bytetrack} during development and both ByteTrack and OATrack
\cite{qi2026oatrack} in the final evaluation. A large fixed input is a common
choice for small aerial objects when no adaptation is used.

\emph{Data and protocol.} Eight sequences of the VisDrone2019-MOT training split
(3282 frames) were declared as calibration data. They are disjoint from the
seven evaluation sequences. All development used four folds, each holding out
two calibration sequences, and reported the cross-validated (CV) mean change
in HOTA against the baseline. Detector outputs were cached for every
resolution and NMS value; a score threshold applied to the cache was checked
to reproduce direct inference exactly. Table~\ref{tab:devrecord} summarizes
the development steps.

\emph{Cue audit.} Stages 1 and 2 searched the state-to-resolution and
state-to-NMS maps on exhaustive grids. Stage 3A searched all 31 non-empty
subsets of the five cues and selected \{crowd\}, which defines the
crowd-selected SCI of Eq.~\eqref{eq:sci}. The other cues were kept as guard and
auxiliary cues of the state logic (Sect.~\ref{sec:method}).

\emph{Score threshold.} At 1088 pixels and NMS 0.45, a fixed score threshold
of 0.40 raised pooled calibration HOTA by $+0.881$ over 0.01 and improved
seven of the eight sequences. The score threshold was therefore added as the
third action dimension.

\emph{Final joint search.} The final search used 60 TPE trials (seed 43). It
searched $\tau_m\in[0.10,0.50]$, $\tau_h\in[\tau_m+0.05,\min(\tau_m+0.40,0.90)]$,
and for each state a resolution in $\{1088,1280,1536\}$, an NMS value in
$\{0.45,0.60,0.70\}$, and a score threshold in $\{0.10,0.25,0.40\}$. The
objective was the CV mean change in HOTA against the baseline, with ByteTrack
as host. Eligible candidates had to realize at least two distinct actions,
each on at least 10\% of calibration frames, and switch within sequences. The
frozen policy is the eligible candidate whose three states use three
different actions (Table~\ref{tab:actions}). On calibration data it used LOW,
MEDIUM, and HIGH on 7.6\%, 10.1\%, and 82.3\% of frames, with 0.607 state
changes per 100 frames on average.

\emph{Comparison.} Table~\ref{tab:folds} compares the frozen policy with the
baseline on each calibration fold. All four folds are positive, with a CV
mean of $+2.447$ HOTA. After these steps the thresholds, cue subset, and
action map were frozen. The final comparison on the evaluation sequences is
given in Sect.~\ref{sec:results}.

\begin{table}[t]
\caption{Main modern development steps on the eight calibration sequences
(ByteTrack host). The score-threshold gain is pooled HOTA; the joint-search
result is the mean over four folds.}
\label{tab:devrecord}
\centering\footnotesize
\setlength{\tabcolsep}{3pt}
\begin{tabular}{@{}lll@{}}
\toprule
Step & Search & Result\\
\midrule
Cue subset & 31 subsets & \{crowd\}\\
Score threshold & 4 values & 0.40; $+0.881$ HOTA\\
Joint search & 60 TPE trials & CV mean $+2.447$\\
\bottomrule
\end{tabular}
\end{table}

\begin{table}[t]
\caption{Frozen policy versus the modern baseline (YOLO11m at 1536 pixels and
NMS 0.70, ByteTrack) on the four calibration folds: change in HOTA.}
\label{tab:folds}
\centering\footnotesize
\begin{tabular}{@{}lrrrrr@{}}
\toprule
 & Fold 1 & Fold 2 & Fold 3 & Fold 4 & Mean\\
\midrule
$\Delta$HOTA & $+0.592$ & $+4.164$ & $+2.029$ & $+3.004$ & $+2.447$\\
\bottomrule
\end{tabular}
\end{table}
